\chapter{Normas IRAM y Escalas}

\section{Normalización IRAM y Aplicación de Escalas}

El dibujo técnico no sería posible si no estuviera estandarizado. Las
normas IRAM (Instituto Argentino de Normalización y Certificación)
garantizan que el diseño de un componente sea interpretado de la misma
manera por ingenieros en Argentina y en otros países que siguen normas
ISO equivalentes.

\subsection{Importancia de la Normalización}

La normalización establece reglas comunes para:
\begin{itemize}
	\item Formatos de papel (A4, A3, A2, etc.).
	\item Tipos de líneas (continua, trazo, punto).
	\item Representación de vistas (proyecciones).
	\item Acotación.
	\item Simbología técnica.
\end{itemize}

Sin estas reglas, un plano de una máquina simple sería incomprensible
para cualquier persona que no haya participado en su diseño.

\subsection{El concepto de Escala}

El uso de escalas permite representar objetos que son o muy grandes o
muy pequeños en un formato de papel estándar. La fórmula básica para la
escala es:

$Escala = \frac{\text{Tamaño Dibujo}}{\text{Tamaño Real}}$

Donde:
\begin{itemize}
	\item \textbf{Tamaño Dibujo}: Medida representada en el papel
	      (en mm).
	\item \textbf{Tamaño Real}: Medida real del objeto (en mm).
\end{itemize}

Existen tres tipos principales de escalas:
\begin{enumerate}
	\item \textbf{Escala de ampliación (E > 1):} El dibujo es más
	      grande que el objeto real (ej: 2:1, 5:1). Se utiliza para
	      objetos muy pequeños.
	\item \textbf{Escala natural (E = 1):} El dibujo tiene el mismo
	      tamaño que el objeto real (1:1).
	\item \textbf{Escala de reducción (E < 1):} El dibujo es más
	      pequeño que el objeto real (ej: 1:2, 1:5, 1:10). Se utiliza
	      para piezas grandes o conjuntos.
\end{enumerate}

Es crucial indicar siempre la escala utilizada en el rótulo del plano.

\section{Cuestionario}
\begin{enumerate}
	\item ¿Qué entidad regula las normas IRAM?
	\item Define qué es la escala natural.
	\item ¿Por qué se usan escalas de reducción?
	\item Menciona tres tipos de líneas normalizadas y su aplicación.
	\item ¿Cómo influye el tamaño del formato (A4, A3, etc.) en la elección de la escala?
	\item Explica la diferencia entre escala gráfica y escala numérica.
	\item ¿Por qué las vistas ortogonales son fundamentales en el dibujo técnico?
	\item ¿Cuál es la finalidad de un rótulo en un plano?
\end{enumerate}

\section{Parte Práctica}
\begin{enumerate}
	\item Diseñar un formato A4 con cajetín (rótulo) según norma IRAM.
	\item Dibujar un rectángulo de 200x150 mm en escala 1:2.
	\item Representar un objeto de 500x300 mm en escala 1:5.
	\item Dibujar una línea de eje usando el estilo de línea adecuado.
	\item Crear una escala 2:1 para un pequeño tornillo de 10mm.
	\item Insertar texto con formato normalizado.
	\item Dibujar una pieza simple aplicando tres escalas distintas.
	\item Completar el rótulo del formato con datos ficticios.
	\item Crear un estilo de cota ajustado a la escala del dibujo.
	\item Realizar un plano de conjunto a escala 1:10.
\end{enumerate}
